\section{Core 1.3 Foundational Consistency Dossier}
\label{sec:oc13-foundational-consistency}

This chapter fixes the internal consistency contract of the flagship volume.
Its purpose is not to replace the formal chapters with editorial commentary,
but to make three things readable in one place:

\begin{itemize}
\item why the present hierarchy is fixed as \(K_0\) through \(K_{12}\) rather
      than by an arbitrary cutoff;
\item the full statement-class discipline for axiomatic, derived, empirical,
      structural, interpretive, and frontier claims;
\item how the master monograph and the narrower journal-core extraction are
      related without silent drift in scope or notation.
\end{itemize}

\subsection{Why the hierarchy extends through \texorpdfstring{$K_{12}$}{K12}}

Core 1.2 already contains enough structure to motivate a release hierarchy
larger than \(K_{10}\). The point is not that higher labels are rhetorically
attractive. The point is that the audited source corpus already distinguishes:

\begin{enumerate}
\item lower-level continua whose axes, thresholds, and flows are internal to a
      given model class;
\item theoretical and meta-theoretical continua that organise theories and
      model-spaces themselves;
\item a further stratum in which operators, meta-logics, and model-spaces
      become explicit objects of structural evolution;
\item and an upper bounding layer in which global coherence conditions across
      such higher-order structures must themselves be represented.
\end{enumerate}

That is why Core 1.3 carries the public hierarchy through \(K_{12}\) rather
than stopping at \(K_{10}\). This is a release-level design rationale, not a
standalone impossibility proof. The supporting route is the K-level doctrine in
Section~\ref{sec:klevels-full}, the theorem-navigation discussion in
Section~\ref{sec:oc13-theorem-roadmap}, and the worked preservation test in
Section~\ref{sec:oc13-worked-examples}. Those sections assign \(K_{11}\) to
evolving meta-theoretical objects and \(K_{12}\) to global compatibility and
integrability conditions. The present release does not claim a proof that no
further extension is conceivable in principle; it claims only that \(K_{12}\)
is the highest level lawfully carried by the current audited source corpus and
therefore the last public level of this release.

\subsection{Scientific role of the upper levels}

\paragraph{Level \texorpdfstring{$K_{11}$}{K11}.}
\(K_{11}\) is the first level at which formal ontologies, operator systems,
meta-logics, and model-spaces are themselves treated as structured continua
subject to admissibility, transformation, and collapse conditions.
This level is needed once one studies not only theories, but the evolution of
the spaces in which theories and theory-building operators live.

\paragraph{Level \texorpdfstring{$K_{12}$}{K12}.}
\(K_{12}\) is not introduced as a mystical completion layer.
It is the minimally specified upper coherence stratum required when one needs
to discuss global compatibility, boundedness, and failure modes across families
of \(K_{11}\)-objects.
Its role is structural rather than triumphalist: it prevents the hierarchy from
pretending that there is no formal problem above recursively evolving
meta-theoretical spaces.

\subsection{Statement classes}

To keep the monograph readable without blurring scientific status, Core~1.3 uses the
following statement discipline throughout the release:

\begin{itemize}
\item \textbf{Axiomatic statements} fix primitive assumptions, admissibility
      conditions, or identity rules.
\item \textbf{Derived statements} are theorem-bearing statements with explicit support
      dependence and proof obligations.
\item \textbf{Empirical statements} require a standalone source-owned domain
      packet, pinned data route, numerical parameter packet, replay procedure,
      and explicit falsifier before promotion.
\item \textbf{Structural statements} describe the organisation of the framework,
      the hierarchy, and the mapping between modules or levels.
\item \textbf{Interpretive statements} explain what the formal structure means,
      which reader it is for, and what later projections may test.
\item \textbf{Frontier statements} may remain in the research program, but may
      not enter positive outward science without formal promotion.
\end{itemize}

These prose classes are a local proof-status taxonomy, not a replacement for
the canonical SPOT verdict vocabulary. The mapping is:
\begin{itemize}
\item Axiomatic and derived statements map to \texttt{THEOREM\_NATIVE} when
      the support route is closed.
\item Empirical statements map to \texttt{THEOREM\_NATIVE} only after the
      packet, data route, numerical packet, replay procedure, and falsifier
      gates pass.
\item Structural and interpretive statements map to \texttt{BRIDGE\_ONLY},
      \texttt{FRAME\_ONLY}, or \texttt{EXTENSION} as recorded by SPOT.
\item Frontier statements map to \texttt{EXTENSION} unless explicitly refuted
      or promoted.
\item Blocked or refuted statements map to \texttt{REFUTED} when SPOT records
      closure failure.
\end{itemize}
This chapter classifies the statement classes that affect theorem status,
structural scope, interpretation, empirical promotion, and frontier boundaries.
Operational route discipline is handled later in the theorem-roadmap and
proof-machinery chapters.
Editorial-bridge language is not a seventh proof class; it is an interpretive
and release-consistency label for statements that connect the monograph,
journal core, dossiers, and public package without adding theorem force.
Operational statements are the later exact-match, held-out replay,
five-sigma/tail-bar, comparator, falsifier, and escalation rules attached to
empirical or structural classes; they are handled after the statement class
has already been fixed here and are not a seventh proof-status class.

The journal-core extraction is allowed to narrow this body, but not to blur
these classes into one another.
An interpretive sentence may explain a theorem; it may not replace the theorem.
A structural remark may motivate a projection; it may not silently serve as
proof of that projection.

\subsection{Proof-bearing and nonclaim layers}

The master monograph therefore separates three layers explicitly:

\begin{enumerate}
\item the inherited formal backbone that still bears the structure of the core
      theory;
\item the Core 1.3 repair layer that clarifies witness discipline, theorem
      fate, chronology, and publication boundary;
\item the bounded nonclaim layer that states what the present release does not
      yet defend, even when the broader research program still pursues it.
\end{enumerate}

This separation matters for hostile review.
A reviewer should be able to ask, for any sentence that sounds strong:
is it axiomatic, derived, empirical, structural, interpretive, or frontier?
If it is derived, what supports it?
If it is empirical, what packet, data route, replay, comparator, and falsifier
support it?
If it is structural or interpretive, what exactly is it not claiming?
If it is frontier, where is it visibly barred from positive release authority?

\subsection{Relation between master monograph and journal core}

The canonical SPOT is the scientific authority of the release, and the
source-owned science corpus is the authoring and witness locus behind that
projection. The master monograph is the flagship reader-facing reference
because it carries the full theory, the reader ladder, the appendices, and the
explicit publication boundary in one place.
The \emph{Journal Core} is the bounded article artifact. The
\emph{journal-core extraction} is the operation that selects one defensible
theorem-bearing route for editorial inspection without pretending that the
whole monograph is already compressed into that shorter artifact.

For that reason the extraction rule is one-way:

\[
\text{Master Monograph} \longrightarrow \text{Journal Core},
\]

with explicit preservation of notation, nonclaim boundaries, and theorem scope.
The extraction procedure may omit detail. The Journal Core may not enlarge the
claim.
The public closure boundary remains the
\texttt{CORE\_1\_3\_SCIENCE\_ONLY} scope stated in
Section~\ref{sec:oc13-publication-boundary}. \texttt{FRAME\_ONLY} closures
may be \texttt{PASS} as frame support, but they do not become standalone
empirical promotions. \texttt{EXTENSION} and \texttt{REFUTED} lanes remain
outside public \texttt{PASS} authority.

\subsection{Reader-facing consequence}

This chapter exists because scientific rigor and readability do not pull in
opposite directions here.
The more explicit the dossier becomes about hierarchy, statement class, and
proof boundary, the easier it becomes for three readers at once:

\begin{itemize}
\item the editor who needs a clean overview of what is actually defended;
\item the formal reader who needs exact scope and dependency discipline;
\item the non-specialist scientific reader who needs a coherent narrative entry
      point before diving into the full technical shell.
\end{itemize}

\paragraph{Bridge to the next chapter.}
With the statement classes, hierarchy boundary, and master-to-derivative
relation fixed, the next chapter narrows further to one task only: the shortest
lawful proof-navigation route through the load-bearing theorem spine.

\clearpage
